\scanpage{166}
% L1–10
On peut construire un torseur
$Q_{\sigma}\deleted{\unsure{(L)}}=Q(L,\sigma)$ sous $L^{\otimes 2}$, et une fonction
\unsure{para-quadratique|sous-quadratique}
\[
A\longrightarrow Q
\]
dont la partie homogène soit
\[
\unsure{x|u}\,\unsure{\vdash|\mapsto}\,u^{\otimes 2}:L\mathbin{\text{---}}L^{\otimes 2},
\]
et qui soit \uline{universelle} pour ces
propriétés, \deleted{Pour ceci, on note que}
\inserted{de sorte que les \deleted{fonctions}}
\unsure{ft-quadratiques|f.-quadratiques|sous-quadratiques} à val. dans $L^{\otimes 2}$ comp. avec
$\sigma$ soient simplement les composées
\diagram{\begin{tikzcd}[ampersand replacement=\&,column sep=large]
 A\arrow[r]\& Q\arrow[r, "{\substack{\mathrm{iso.}\ \mathrm{de}\\L^{\otimes 2}\text{-tors.}\\(\text{\unsure{correspond$^{\mathrm t}$|correspondant}}\\\text{au choix d'une}\\\text{origine})}}"']\&L^{\otimes 2}
\end{tikzcd}}

% L12–18
Pour ceci on \unsure{note|voit} que l'ens des
fonctions \deleted{\unsure{sur|sous}} quadr. (à valeurs dans $L^{\otimes 2}$)
qui nous conviennent, forment un
torseur sous le module des formes
quadr. qui se factorisent par
$E/\unsure{E^0|E_0|E^\circ}\simeq k$, lequel module est \unsure{isom|iso}
à $\mathrm{Quad}(k,L^{\otimes 2})\simeq L^{\otimes 2}$. En termes
% archive sheet number: 165
